% GENERATED FILE. Edit canonical lessons, then run npm run generate:books.
% canonical-source-hash: fnv1a64:2df88d26b1bceb54
% canonical-lessons: MW-C180-dharno, MW-C180-ghalno, MW-C180-nakhno, MW-C180-chhiko, MW-C180-chhabri

\chapter{Putting Things Away}
\label{ch:mw-putting-things-away}
\hlchaptermodality{180}

\begin{chapteropening}
\textbf{By the end of this chapter:} \emph{I can say where I put things and what I throw out.}

Complete the last lesson of chapter 180: I can say where I put things and what I throw out.
\end{chapteropening}

\section[\texorpdfstring{\mw{धरणो} (dharṇo)}{धरणो (dharṇo)}]{\mw{धरणो} (dharṇo) — to set down; to place}

\label{lesson:MW-C180-dharno}



\begin{quote}
Before the new one: say the Marwadi for to peel, then the Marwadi for to winnow.
\end{quote}

\subsection*{You'll want to know: \mw{धरणो}}
\textbf{\mw{धरणो}} — \emph{dharṇo} — \textquotedblleft{}to set down; to place\textquotedblright{}.

\begin{quote}
\mw{हूं} \mw{घड़ो} \mw{भींत} \mw{कनै} \mw{धरूं} \mw{हूं।}
\end{quote}

\begin{quote}
\emph{hū̃ ghaṛo bhī̃t kanai dharū̃ hū̃.}
\end{quote}

\begin{quote}
I set the water pot down by the wall.
\end{quote}

Past (I …): \mw{म्हैं} \mw{घड़ो} \mw{भींत} \mw{कनै} \mw{धर्यो}, \emph{mhaĩ ghaṛo bhī̃t kanai dharyo}.

\begin{cousinweb}
\textbf{\mw{धरणो}} (\emph{dharṇo}) is inherited from Sanskrit \emph{dharati}, \textquotedblleft{}holds, keeps\textquotedblright{}; CDIAL lists Marwari \emph{dharṇo}, \textquotedblleft{}to keep, place\textquotedblright{} (CDIAL 6747).
\end{cousinweb}

\begin{grammarlens}[title={What you've built}]
One more word for this chapter.
\end{grammarlens}

\subsection*{Your turn}
\begin{itemize}
\raggedright
  \item \emph{Say it:} \emph{dharṇo}
  \item \emph{Say it:} \emph{dharṇo}, once more
  \item \emph{Say it:} \emph{hū̃ ghaṛo bhī̃t kanai dharū̃ hū̃.} — I set the water pot down by the wall.
  \item \emph{Recall:} say the Marwadi for to peel, then the Marwadi for to winnow, then say \emph{dharṇo} again
\end{itemize}

\begin{tcolorbox}[breakable,colback=teal!4,colframe=teal!35!black,title={Before you move on}]
What does \mw{धरणो} mean? (\textquotedblleft{}To set down; to place\textquotedblright{}.) Say it once more.
\end{tcolorbox}
\section[\texorpdfstring{\mw{घालणो} (ghālṇo)}{घालणो (ghālṇo)}]{\mw{घालणो} (ghālṇo) — to put in}

\label{lesson:MW-C180-ghalno}



\begin{quote}
Before the new one: say the Marwadi for to winnow, then the Marwadi for to set down.
\end{quote}

\subsection*{You'll want to know: \mw{घालणो}}
\textbf{\mw{घालणो}} — \emph{ghālṇo} — \textquotedblleft{}to put in\textquotedblright{}.

\begin{quote}
\mw{मां} \mw{दाल} \mw{मांय} \mw{लूण} \mw{घालै} \mw{है।}
\end{quote}

\begin{quote}
\emph{mā̃ dāl mā̃y lūṇ ghālai hai.}
\end{quote}

\begin{quote}
Mother puts salt in the dal.
\end{quote}

Past (I …): \mw{म्हैं} \mw{दाल} \mw{मांय} \mw{लूण} \mw{घाल्यो}, \emph{mhaĩ dāl mā̃y lūṇ ghālyo}.

\begin{cousinweb}
\textbf{\mw{घालणो}} (\emph{ghālṇo}) is inherited from a reconstructed \emph{*ghalyati}, \textquotedblleft{}pours\textquotedblright{}; CDIAL lists Old Marwari \emph{ghālaï}, \textquotedblleft{}puts, places, gives\textquotedblright{} (CDIAL 4457).
\end{cousinweb}

\begin{grammarlens}[title={What you've built}]
One more word for this chapter.
\end{grammarlens}

\subsection*{Your turn}
\begin{itemize}
\raggedright
  \item \emph{Say it:} \emph{ghālṇo}
  \item \emph{Say it:} \emph{ghālṇo}, once more
  \item \emph{Say it:} \emph{mā̃ dāl mā̃y lūṇ ghālai hai.} — Mother puts salt in the dal.
  \item \emph{Recall:} say the Marwadi for to winnow, then the Marwadi for to set down, then say \emph{ghālṇo} again
\end{itemize}

\begin{tcolorbox}[breakable,colback=teal!4,colframe=teal!35!black,title={Before you move on}]
What does \mw{घालणो} mean? (\textquotedblleft{}To put in\textquotedblright{}.) Say it once more.
\end{tcolorbox}
\section[\texorpdfstring{\mw{नाखणो} (nākhṇo)}{नाखणो (nākhṇo)}]{\mw{नाखणो} (nākhṇo) — to throw (away)}

\label{lesson:MW-C180-nakhno}



\begin{quote}
Before the new one: say the Marwadi for to set down, then the Marwadi for to put in.
\end{quote}

\subsection*{You'll want to know: \mw{नाखणो}}
\textbf{\mw{नाखणो}} — \emph{nākhṇo} — \textquotedblleft{}to throw (away)\textquotedblright{}.

\begin{quote}
\mw{हूं} \mw{रोज} \mw{कचरो} \mw{नाखूं} \mw{हूं।}
\end{quote}

\begin{quote}
\emph{hū̃ roj kachro nākhū̃ hū̃.}
\end{quote}

\begin{quote}
I throw out the rubbish every day.
\end{quote}

Past (I …): \mw{म्हैं} \mw{कचरो} \mw{नाख्यो}, \emph{mhaĩ kachro nākhyo}.

\begin{cousinweb}
\textbf{\mw{नाखणो}} (\emph{nākhṇo}) is inherited, by CDIAL's account, from Sanskrit \emph{naṅkṣyati}, \textquotedblleft{}will perish\textquotedblright{}, a sense far from today's; CDIAL lists Old Marwari \emph{nākhaï}, \textquotedblleft{}puts, drops, strews\textquotedblright{} (CDIAL 6930).
\end{cousinweb}

\begin{grammarlens}[title={What you've built}]
One more word for this chapter.
\end{grammarlens}

\subsection*{Your turn}
\begin{itemize}
\raggedright
  \item \emph{Say it:} \emph{nākhṇo}
  \item \emph{Say it:} \emph{nākhṇo}, once more
  \item \emph{Say it:} \emph{hū̃ roj kachro nākhū̃ hū̃.} — I throw out the rubbish every day.
  \item \emph{Recall:} say the Marwadi for to set down, then the Marwadi for to put in, then say \emph{nākhṇo} again
\end{itemize}

\begin{tcolorbox}[breakable,colback=teal!4,colframe=teal!35!black,title={Before you move on}]
What does \mw{नाखणो} mean? (\textquotedblleft{}To throw (away)\textquotedblright{}.) Say it once more.
\end{tcolorbox}
\section[\texorpdfstring{\mw{छींको} (chhī̃ko)}{छींको (chhī̃ko)}]{\mw{छींको} (chhī̃ko) — a hanging net for pots}

\label{lesson:MW-C180-chhiko}



\begin{quote}
Before the new one: say the Marwadi for to put in, then the Marwadi for to throw.
\end{quote}

\subsection*{You'll want to know: \mw{छींको}}
\textbf{\mw{छींको}} — \emph{chhī̃ko} — \textquotedblleft{}a hanging net for pots\textquotedblright{}.

\begin{quote}
\mw{छींको} \mw{नुवो} \mw{है।}
\end{quote}

\begin{quote}
\emph{chhī̃ko nuvo hai.}
\end{quote}

\begin{quote}
The hanging net is new.
\end{quote}

A masculine noun.

\begin{cousinweb}
\textbf{\mw{छींको}} (\emph{chhī̃ko}) is inherited from Sanskrit \emph{śikya}, \textquotedblleft{}a rope sling for carrying things\textquotedblright{}; CDIAL lists Hindi \emph{chī̃kā} and Sindhi \emph{chiko}, a net or shelf that hangs (CDIAL 12427). CDIAL gives no Marwari form in this entry.
\end{cousinweb}

\begin{grammarlens}[title={What you've built}]
One more word for this chapter.
\end{grammarlens}

\subsection*{Your turn}
\begin{itemize}
\raggedright
  \item \emph{Say it:} \emph{chhī̃ko}
  \item \emph{Say it:} \emph{chhī̃ko}, once more
  \item \emph{Say it:} \emph{chhī̃ko nuvo hai.} — The hanging net is new.
  \item \emph{Recall:} say the Marwadi for to put in, then the Marwadi for to throw, then say \emph{chhī̃ko} again
\end{itemize}

\begin{tcolorbox}[breakable,colback=teal!4,colframe=teal!35!black,title={Before you move on}]
What does \mw{छींको} mean? (\textquotedblleft{}A hanging net for pots\textquotedblright{}.) Say it once more.
\end{tcolorbox}
\section[\texorpdfstring{\mw{छाबड़ी} (chhābṛī)}{छाबड़ी (chhābṛī)}]{\mw{छाबड़ी} (chhābṛī) — a flat basket}

\label{lesson:MW-C180-chhabri}



\begin{quote}
Before the new one: say the Marwadi for to throw, then the Marwadi for a hanging net.
\end{quote}

\subsection*{You'll want to know: \mw{छाबड़ी}}
\textbf{\mw{छाबड़ी}} — \emph{chhābṛī} — \textquotedblleft{}a flat basket\textquotedblright{}.

\begin{quote}
\mw{छाबड़ी} \mw{मांय} \mw{रोटी} \mw{है।}
\end{quote}

\begin{quote}
\emph{chhābṛī mā̃y roṭī hai.}
\end{quote}

\begin{quote}
There are rotis in the flat basket.
\end{quote}

A feminine noun.

\begin{cousinweb}
\textbf{\mw{छाबड़ी}} (\emph{chhābṛī}) is inherited, by CDIAL's account, from Sanskrit \emph{chadman}, \textquotedblleft{}a roof\textquotedblright{}; CDIAL lists Hindi \emph{chabṛī}, \textquotedblleft{}a strainer, a bamboo basket\textquotedblright{} (CDIAL 4981). CDIAL gives no Marwari form in this entry.
\end{cousinweb}

\begin{grammarlens}[title={What you've built}]
One more word for this chapter.
\end{grammarlens}

\subsection*{Your turn}
\begin{itemize}
\raggedright
  \item \emph{Say it:} \emph{chhābṛī}
  \item \emph{Say it:} \emph{chhābṛī}, once more
  \item \emph{Say it:} \emph{chhābṛī mā̃y roṭī hai.} — There are rotis in the flat basket.
  \item \emph{Recall:} say the Marwadi for to throw, then the Marwadi for a hanging net, then say \emph{chhābṛī} again
\end{itemize}

\begin{tcolorbox}[breakable,colback=teal!4,colframe=teal!35!black,title={Before you move on}]
What does \mw{छाबड़ी} mean? (\textquotedblleft{}A flat basket\textquotedblright{}.) Say it once more.
\end{tcolorbox}
